\begin{claim}{D.2.33}{Proved}{Every once-suspended map has the same remaining Toda test}
Keep the spaces and the single normalization of
\textup{D.2.27--D.2.28}, and let
$\lambda\in\pi_8W\cong\ZZ/2$ be the class whose mod-two
Hurewicz image is $R_8$. The Toda bracket
\[
 \Delta=\langle\lambda,\eta_8,2\rangle
       \subseteq\pi_{10}W
 \tag{D.2.33a}
\]
is a singleton, and its value is either $0$ or the element
$p$ of \textup{D.2.32}. If an ordinary map
$h:\Sigma A\to Y$ has the nonzero coefficient
$\epsilon=1$ of \textup{D.2.28}, then
\[
                         \Delta=p.
 \tag{D.2.33b}
\]
More precisely, its quadratic stable summand has a definite
Moore-top choice $\tau\in\FF_2$, relative to the reference
Moore map fixed in \textup{D.2.32}; its eleven-cell
obstruction, modulo the integral-Hurewicz kernel
$\langle s\rangle=\langle\rho\nu\rangle$, is
\[
                    p+\tau\Delta.
 \tag{D.2.33c}
\]
Thus $\epsilon=1$ requires both $\Delta=p$ and $\tau=1$.
In particular, $\Delta=0$ would imply that the normalized
functional is zero for every once-suspended attachment.
The value of this bracket is determined in \textup{D.2.34} below;
the present necessary condition does not construct a map on the full
source or prove a statement about subsequent nullification stages.
\end{claim}
\begin{proof}
Write $X=\Sigma B$, $A=QX$, and
$Z=\fib(Y\to K(\FF_2,4))$, so $W=\Omega Z$.
All cell calculations are two-local. The homotopy groups used
are the actual finite two-primary groups already computed.

\medskip\noindent
\textbf{Exactly the stabilization needed for an obstruction.}
The space $Z$ is five-connected, with
$\pi_6Z=\ZZ/2\{\chi_Z\}$. The double adjunct of its
Whitehead square is the Browder square in \textup{D.2.23},
so $[\chi_Z,\chi_Z]=0$ in $\pi_{11}Z$.
Use the reduced-product model $JZ\simeq\Omega\Sigma Z$
\cite[Section~1 and Theorem~5.6]{I_James1955}.
The pair $(JZ,Z)$ is eleven-connected. Its first relative
group is
\[
 \pi_{12}(JZ,Z)\cong H_{12}(JZ,Z)
       \cong H_6Z\otimes H_6Z\cong\ZZ/2.
 \tag{D.2.33d}
\]
This follows directly from relative Hurewicz and the
word-length filtration: two letters first occur in degree
twelve, and three letters cannot affect that group.
Naturality from $S^6\to Z$ identifies its boundary with
$[\chi_Z,\chi_Z]$, since the twelve-cell of $J_2S^6$
attaches by the ordinary Whitehead square. The exact sequence
therefore makes $\pi_{11}Z\to\pi_{12}\Sigma Z$ an
isomorphism. The latter group is already stable because
$\Sigma Z$ is six-connected. Hence
\[
 \pi_{11}Z\xrightarrow{\ \cong\ }
       \pi_{11}\Sigma^\infty Z.
 \tag{D.2.33e}
\]
In particular $W\to\Omega QZ$ is ten-connected and
injective on $\pi_{10}$. This is a statement about the
specified obstruction group, not an equivalence of
unrestricted mapping sets.

Suppose now that a full map $g:A\to W$ has
$g^*\vartheta_7=\alpha_3\beta_4$. Stabilize its adjunct
$\Sigma A\to Z$ and use the natural Snaith splitting
\cite[Theorem~1.1 and Section~2]{Snaith1974}. Its quadratic
summand gives a map
\[
       D_2X=(E\Sigma_2)_+\wedge_{\Sigma_2}X^{\wedge2}
                       \longrightarrow\Omega QZ.
 \tag{D.2.33f}
\]
The actual space $D_2X$ is five-connected. The ten-connectivity
just proved lifts its ten-skeleton to $W$. Each eleven-cell
obstruction has zero image in $\pi_{10}\Omega QZ$, since
\textup{(D.2.33f)} is defined on the full summand.
Injectivity in \textup{(D.2.33e)} makes every actual obstruction
zero. We have therefore obtained an actual map on the eleven-stage
of $D_2X$ preserving the quadratic $\vartheta_7$ functional.
Here $\vartheta_7$ is the suspension of the ordinary class
$\vartheta_8$ on $Z$, so its pairing is preserved by this
stabilization and adjunction. No lifting of an arbitrary
twelve-dimensional stable map is being assumed.

This also separates all the unit lower data from the present
test. In the stable Snaith grading the classes
$AX,BD,BE,BX$ have weight two. The classes on the unit
$X$, including possible nonzero data on its five- and
six-cells, have weight one. They cannot contribute to the
restriction \textup{(D.2.33f)}.

\medskip\noindent
\textbf{The complete quadratic nine-stage.}
Use a finite two-local cellular model of this summand through
dimension eleven, with odd-unit pairs cancelled. Its attaching
groups through dimension ten are in the ordinary stable range,
since its bottom dimension is six. Abbreviate the unit homology
labels as in \textup{D.2.32}, and put
$R=Q^4A$, $P=AC$, $N=Q^5A$, $J=AB+R$.
The weight-two part of the previously computed eight-stage is
\[
 C\vee M(\ZZ/4,7)_J\vee S^8_s,
 \qquad
 C=C_1^8=\cofib((2,\eta_6):S^6\vee S^7\to S^6),
 \qquad s=N-P.
 \tag{D.2.33g}
\]
Its top Moore cell is $B^2-P$, while $P$ is the top cell
of $C$. Write
$\pi_8C=\ZZ_{(2)}\{\gamma\}\oplus\ZZ/2\{\delta\}$,
with quotient coordinates $\gamma\mapsto(0,2)$ and
$\delta\mapsto(\eta_7,0)$ as in \textup{D.2.18}.
The five nine-attachments are exactly
\[
\begin{array}{c|ccccc}
 &AD&AE&BC&Q^6A&Q^5B\\\hline
 &\tau_J+\delta&0&\gamma+\tau_J+\delta&
                    \gamma+2s&\delta.
\end{array}
 \tag{D.2.33h}
\]
Integral degrees and the second squares detect every component.
The even-$\nu$ components retained in the full-source table
belong to its unit summand and do not occur here. Cell slides
give the four factors
\[
\begin{split}
 J_*&=M(\ZZ/4,7)_J\cup_{i\eta}e^9_{AD-Q^5B},\\
 C_*&=C\cup_\delta e^9_{Q^5B}
                    \cup_\gamma e^9_{BC-AD},\\
 M(\ZZ/2,8)_s,&\qquad S^9_{AE}.
\end{split}
 \tag{D.2.33i}
\]
All these comparisons are ordinary in this connectivity and
dimension range.

The four integral degree-ten cycle coordinates in this
summand reduce to
\[
             C^2,\quad Q^7A,\quad AF,\quad
                    K=AX+BD+BE.
 \tag{D.2.33j}
\]
The complete Bockstein table in \textup{D.2.32}, restricted
to weight two, gives
$H_{10}(D_2X;\ZZ_{(2)})\cong(\ZZ/2)^4$.
Indeed the mod-two dimensions in degrees nine and ten are
five and seven, the integral eighth group has two cyclic
summands, and the degree-eleven first Bockstein has the four
independent images in \textup{(D.2.33j)}.
For the eighth group, the established weight-two seventh
group is $\ZZ/4$, the mod-two eighth dimension is three,
and the two independent first-Bockstein images from degree
nine fix both exponents as two.
The integral eleven-boundary lattice is twice the cycle lattice.
The exact Pontryagin chain calculation retains
\[
                 d(BX)=2K.
 \tag{D.2.33k}
\]
The cycle $K$ has zero integral ninth-cell degrees and
\[
        \operatorname{Sq}_*^2K=P+B^2,
                  \qquad \operatorname{Sq}_*^4K=0.
 \tag{D.2.33l}
\]

Its $J_*$ projection is precisely the Moore eta-coextension,
by the complete Chang calculation of \textup{D.2.32}.
Its projections to the other three factors vanish.
For $M(\ZZ/2,8)_s$ this follows from the zero corresponding
second-square coefficient, and for the nine-sphere it follows
from the zero integral degree. Here is the full check for $C_*$,
so the nonzero bottom-six branch is not omitted.
The first three stable stems and the Chang cofiber sequence give
\[
 \pi_9C=(\ZZ/2)\{i\nu_6,\delta\eta_8\},
 \qquad \gamma\eta_8=0.
\]
The last equality is the cone-of-$\eta$ fourth-square argument
already used in \textup{D.2.19}; the only possible residual
bottom $\nu$ has nonzero fourth-square detector.
After the two nine-cells the exact sequence yields
\[
 \pi_9C_*\cong\ZZ_{(2)}\{\kappa\}\oplus\ZZ/2\{i\nu_6\},
               \qquad h_{\ZZ}(\kappa)=2[Q^5B].
 \tag{D.2.33m}
\]
The integral degree of $K$ excludes $\kappa$, and its fourth
square excludes $i\nu_6$. Thus there is no undetected component.

\medskip\noindent
\textbf{The nonzero bottom-six branch is excluded.}
The class $A^2$ is the first Bockstein of $Q^4A$.
Every first Bockstein from $H_7W$ is zero, by the complete
low target table. Thus the bottom six-sphere of the quadratic
summand has zero mod-two Hurewicz after mapping to $W$.
Its value is either $0$ or $2\alpha$ in
$\pi_6W=\ZZ/4\{\alpha\}$.

The other ten-cycle $C^2$ has zero integral ninth-cell degrees and
\[
 \operatorname{Sq}_*^2(C^2)=0,
       \qquad \operatorname{Sq}_*^4(C^2)=A^2.
 \tag{D.2.33p1}
\]
Cartan gives the first equality from $\beta C=0$, and the
second from $\operatorname{Sq}_*^2C=A$.
The complete groups of the four factors in
\textup{(D.2.33i)} now identify its attachment as precisely
the bottom $i\nu_6$ in $C_*$.
There is no hidden eta-squared component on a seven-sphere:
the $J_*$ Chang group has only its free generator and the
Moore coextension detected by the second square, and the
other two factors have only the stated eta and degree
coordinates. Equation \textup{(D.2.33m)} accounts for the
entire $C_*$ component.

Suppose the bottom maps to $2\alpha$. Represent the
nine-stage map in $\Sigma L_8$ as justified in the next
paragraph. The $C^2$ attachment then maps to exactly
$2n=2\alpha\nu$ there. There is also an exact integral
relation, not merely its first Bockstein:
\[
                         d(Q^6C)=\pm2C^2.
 \tag{D.2.33p2}
\]
Choose an integral cycle representing the odd-dimensional
class $C_5$. In the arity-two chain model the ordered
configuration one-cell satisfies
$d(e_1(C\otimes C))=-2e_0(C\otimes C)$, since the
transposition of two degree-five labels has Koszul sign
$-1$. Its mod-two classes are exactly $Q^6C$ and $C^2$.
A parity-preserving cellular normalization retains this
relation. Filling the ten-cell by $2N'$ makes the
corresponding eleven-obstruction have integral Hurewicz
$4N'\ne0$. Any other filling changes the obstruction by
twice an element of $\pi_{10}W$, hence by zero.
The $J$-top toggle has zero coefficient on this cycle,
as its second square is zero. This contradicts the
actual eleven-stage extension obtained from
\textup{(D.2.33f)}. The bottom must therefore map to zero.
It may be made constant by a homotopy and collapsed.
This rules out the additional lower-six branch rather
than identifying it with the reference map.

\medskip\noindent
\textbf{Non-J data do not change the BX p contribution.}
Any map from the five-connected nine-stage in
\textup{(D.2.33i)} to $W$ lifts to its five-connected cover
$W'$, then through the eleven-connected evaluation
$\Sigma E\to W'$. The actual comparison
$\Sigma L_8\to\Sigma E$ is nine-connected in the low
cellular realization range. Thus the nine-stage map can be
represented in $\Sigma L_8$. The image of
$\pi_{10}\Sigma L_8$ in $\pi_{10}W$ is contained in
$\langle s\rangle$, by the complete torsion calculation
\textup{(D.2.22t7)--(D.2.22t16)}. In particular it has no
component on $p$, whose integral Hurewicz is $4N'$.
This factorization also justifies the reference filling
used for \textup{(D.2.33p2)}.

After the bottom six-sphere is collapsed, $C_*$ becomes
the wedge of $S^7_R\cup_\eta e^9$ and $M(\ZZ/2,8)_P$.
Its first factor has only a free degree-two generator in
the relevant attaching group, and its second has only the
bottom eta. Both $K$ coordinates are zero, by its degree
and the cancellation in $P+B^2$. This is the quadratic
part of the complete zero-$K_6$ component check in
\textup{D.2.32}; the explicit calculation
\textup{(D.2.33m)} also proves it before collapsing.
Every other lower-factor contribution therefore remains
inside $\Sigma L_8$ and cannot change the $BX$ coefficient
on $p$. Possible $s$-valued obstructions are not discarded.

The inherited functional has value one on $J=AB+Q^4A$,
and zero on $Q^4A$. Therefore the $J$ bottom maps to an
odd multiple of $\rho$. Fix the corresponding odd multiple
of the reference Moore map in \textup{D.2.32}.
All maps with that bottom differ on the Moore top by
$\pi_8W=\ZZ/2\{\lambda\}$. Denote this choice by $\tau$.
Choices on its nine-cell cannot change the $K$ composite,
because the integral degree on that cell is zero.
For the reference top, the Moore/Toda computation of
\textup{(D.2.32j)} sends the eta-coextension to $2n$ in
$\pi_9\Sigma L_8$; an odd multiple gives the same value.
Fill it with $2N'$. All other components of $K$ are null
already in that lower target by \textup{(D.2.33l)--(D.2.33m)}.
The relation \textup{(D.2.33k)} consequently has integral
Hurewicz $4N'$, hence obstruction $p$ modulo $\langle s\rangle$.
Changes of ten-cell fillings add twice an element of
$\pi_{10}W$, which is zero by \textup{D.2.31}.

\medskip\noindent
\textbf{The only change in that direction is one exact Toda value.}
There is an actual quotient $q_M$ from the nine-stage onto
$S^8$, projecting the $J$-Moore top and killing all other
factors. Its nine-cell attaches into the bottom seven-sphere,
so the quotient is well-defined. The other top choice is the
ordinary loop-space sum with $\lambda q_M$.

Put $C\eta_8=S^8\cup_{\eta_8}e^{10}$. Its exact sequence gives
\[
 \pi_9C\eta_8=0,\qquad
 \pi_{10}C\eta_8=\ZZ_{(2)}\{\Gamma\},
                \qquad h_{\ZZ}(\Gamma)=2[e^{10}].
 \tag{D.2.33n}
\]
The pinch is injective with image $2\ZZ_{(2)}$:
multiplication by $\eta$ maps the first stable stem onto
the second, while the following map reduces degree modulo two.
Compose $q_M$ with $S^8\to C\eta_8$ and extend over all
ten-cells. The second-square table gives odd relative top
degree only on $K$, and even degree on the other three cycle
coordinates. Changes by $\Gamma$ adjust these degrees by two,
so they may be chosen as $1$ and $0$, respectively.
The $BX$ attachment then maps to exactly $\Gamma$;
every zero-degree relation maps to zero by
\textup{(D.2.33n)}. This checks the whole relevant coupling,
without assuming a four-cell subcomplex of the source.

Since $\lambda\eta_8=0$ by adjunction from
\textup{(D.2.20h2)}, extend $\lambda$ to $C\eta_8$.
Its value on $\Gamma$ is exactly the bracket in
\textup{(D.2.33a)}. Its full indeterminacy is
\[
 \lambda\circ\pi_{10}S^8+2\pi_{10}W
       =\langle\lambda\eta_8^2\rangle+2\pi_{10}W=0.
 \tag{D.2.33o}
\]
This proves the asserted singleton and the change
$p\mapsto p+\tau\Delta$.

To determine its possible values, lift the extension
$C\eta_8\to W$ through $\Sigma E$ in the same ordinary
range as above. The degree-two pinch of $\Gamma$ makes
the corresponding integral Hurewicz in $H_9E$ twice a
class. Since $\pi_{10}W$ has exponent two, this is either
$0$ or $4N'$: the other summands of
$H_9E=\ZZ/8\oplus(\ZZ/2)^3$ have zero double.
The actual cone comparison
$C\eta_8\cup_\Gamma e^{11}\to C_\Delta$ has relative
top degree one. Its source has no cohomology in degrees
seven or nine, and its top mod-two class is nonzero.
Consequently the suspended detector
$\operatorname{Sq}^4\vartheta_7+
\operatorname{Sq}^2\lambda_9$ has value zero on this cone.
The exact kernel and detector of \textup{D.2.22--D.2.23}
therefore give $\Delta\in\{0,p\}$, with no $s$ addition.

An actual extension on the quadratic eleven-stage must
have zero $BX$ obstruction. The preceding computation
forces $p+\tau\Delta=0$ modulo $\langle s\rangle$.
Since $p$ is nonzero in that quotient, this requires
$\Delta=p$ and $\tau=1$. Applied to the actual
eleven-stage map obtained from \textup{(D.2.33f)}, it proves
the necessary condition for every full ordinary map,
with exactly the scope stated.
\end{proof}
